\section{Experiments}
\label{sec:results}

\subsection{\textsc{Dream} fails to break the ceiling}
\label{sec:dream-results}

Tab.~\ref{tab:dream} reports DREAM under the three modulation
operators, both protocols (P-train and P-zero), and four anchor
budgets $K \in \{1, 3, 5, 10\}$. Probes use the SOTA recipe: bagged
LR (20 bootstraps $\times$ 5 seeds) + val-tuned ordinal thresholds.

\begin{table}[t]
\caption{\textsc{Dream} results on DAiSEE test. Single-LR results
use the standard probe; the final two rows use the SOTA recipe
(bagged LR $K{=}20{\times}5$~seeds + val-tuned ordinal thresholds).
$\Delta$ = $\kappa_q$(method) $-$ matched-recipe raw baseline. All
variants underperform the baseline at both single-LR and bagged
recipes.}
\label{tab:dream}
\centering
\small
\begin{tabular}{llcccc}
\toprule
Op. & Protocol & $K$ & $\kappa_q$ & 95\% CI & $\Delta$\\
\midrule
\multicolumn{6}{c}{\emph{Single LR + threshold tuning} (baseline = 0.232)}\\
Sub & P-train & 1 & 0.000 & $[.000, .000]$ & $-0.232$\\
Sub & P-train & 3 & 0.090 & $[.043, .134]$ & $-0.142$\\
Sub & P-train & 5 & 0.062 & $[.026, .102]$ & $-0.170$\\
Sub & P-train & 10 & 0.082 & $[.044, .121]$ & $-0.150$\\
Sub & P-zero & 5 & 0.085 & $[.044, .126]$ & $-0.147$\\
Cat & P-train & 1 & 0.146 & $[.106, .184]$ & $-0.086$\\
Cat & P-train & 3 & 0.180 & $[.141, .217]$ & $-0.052$\\
Cat & P-train & 5 & 0.222 & $[.174, .269]$ & $-0.010$\\
Cat & P-train & 10 & 0.206 & $[.159, .254]$ & $-0.026$\\
Cat & P-zero & 5 & 0.204 & $[.154, .255]$ & $-0.028$\\
FiLM & P-train & 5 & 0.184 & $[.143, .228]$ & $-0.048$\\
\midrule
\multicolumn{6}{c}{\emph{Bagged LR + threshold tuning} (baseline = 0.247)}\\
Sub & P-train & 5 & 0.089 & $[.044, .139]$ & $-0.158$\\
Cat & P-train & 5 & 0.218 & $[.174, .260]$ & $-0.029$\\
\midrule
Raw single LR (baseline) & --- & --- & 0.232 & $[.187, .278]$ & ---\\
Raw bagged + thresh (SOTA) & --- & --- & \textbf{0.247} & $[.198, .290]$ & ---\\
\bottomrule
\end{tabular}
\end{table}

Three patterns are clear:

\paragraph{P1. \textsc{Sub} consistently destroys signal.} Even the
most informative anchor budget ($K = 3$ or $K = 10$, P-train) leaves
$\kappa_q$ at or below $0.10$. At $K = 1$, the LR readout collapses
entirely (val C-tuning picks $C = 10^{-3}$, almost no regularization
capacity). This mirrors the failure mode of linear concept erasure
(LEACE, $\kappa_q = 0.04$) and of subject-mean subtraction (TFIC).

\paragraph{P2. \textsc{Cat} preserves signal but cannot exceed
baseline.} The strongest \textsc{Dream} variant (\textsc{Cat},
$K = 5$, P-train) sits at $\kappa_q = 0.222$, within bootstrap CI
of the raw baseline ($0.238$) but with mean $-0.016$ below. The
addition of residual features as an extra block does not introduce
information that LR can exploit; the val-tuned $C$ regularizes the
residual features toward $0$, falling back on the raw signal.

\paragraph{P3. P-zero degrades P-train.} Picking anchors from the
first 10\% of frames in temporal order (rather than from the
lowest-blendshape frames) sits about $0.02$ below P-train, but the
direction of the effect is consistent. The lowest-blendshape
criterion is a marginally better proxy for ``identity reference''
than temporal-prefix.

\subsection{The DREAM failure mode is the diagnosis}
\label{sec:dream-diagnosis}

Tab.~\ref{tab:diag-dream} runs the subject-ID and within-subject
ranking diagnostics on DREAM-modulated features. If \textsc{Dream}
were doing anything useful at all, subject-ID accuracy on
$\phi_{\mathrm{sub}}(x) = f(x) - \mathbf{a}_s$ should drop
significantly below 99.5\%; in practice it drops only to
$\approx 92\%$ (still high) and engagement $\kappa_q$ collapses
to $0.08$. The within-subject Spearman $\rho$ stays at $\approx 0.05$,
confirming that \textsc{Dream}-sub does not provide any new
per-clip discrimination.

\begin{table}[t]
\caption{Subject-ID probe and within-subject Spearman on DREAM
features. ``id-acc'' is within-train-subject 80/20 classification
accuracy of subject identity; ``$\rho$'' is mean Spearman rank
correlation between $\mathbb{E}[y]$ and true engagement within each
test subject (n=21 subjects).}
\label{tab:diag-dream}
\centering
\small
\begin{tabular}{lccc}
\toprule
Features & id-acc & engagement $\kappa_q$ & within-subj $\rho$\\
\midrule
Raw SigLIP-L & 99.5\% & 0.199 & 0.052\\
$\phi_{\mathrm{sub}}$, $K=5$ & 92.4\% & 0.077 & 0.043\\
$\phi_{\mathrm{cat}}$, $K=5$ & 99.6\% & 0.207 & 0.067\\
LEACE-erased & 0.9\% & 0.038 & 0.013\\
\bottomrule
\end{tabular}
\end{table}

The failure mode confirms the entanglement diagnosis: removing
identity-aligned variance from frozen features removes engagement
signal proportionally. \textsc{Dream}'s self-supervised anchor is
the gentlest identity removal we can construct; it is not
gentle enough.

\subsection{Cross-task control: FER2013}
\label{sec:fer2013}

To confirm that the failure is engagement-specific rather than a
defect of the encoder, we evaluate frozen CLIP-B/32 and DINOv2-B/14
features on FER2013~\cite{goodfellow2013fer2013} (7-class facial
emotion, 35,887 images, random 80/20 split):

\begin{itemize}\itemsep0pt
  \item CLIP-B/32 LR: accuracy = 62.8\%, $\kappa = 0.554$
  \item DINOv2-B/14 LR: accuracy = 62.1\%, $\kappa = 0.544$
\end{itemize}

The same encoders that ceiling at $\kappa_q \approx 0.10$ on DAiSEE
engagement reach $\kappa \approx 0.55$ on FER2013 emotion. The
failure is not a property of the encoder per se; it is a property
of the engagement task as represented in DAiSEE, where subject
identity dominates and the ordinal label is correlated with
identity-level priors.

\subsection{Sensitivity to encoder scale}
\label{sec:scale}

Tab.~\ref{tab:encoders} reports the best probe per encoder.

\begin{table}[t]
\caption{Best frozen-feature recipe (bagged LR + threshold) on seven
encoders. The ceiling shifts upward modestly with image-encoder scale
($\sim 0.04$ between CLIP-B/32 and SigLIP-L) but plateaus at
SigLIP-L; the larger SigLIP-SO400M does \emph{not} continue the
trend. \textbf{Video-temporal pretraining underperforms image VLMs}:
VideoMAE-base (pre-trained on Kinetics action recognition with 16
frames per clip) achieves only $\kappa_q = 0.101$.}
\label{tab:encoders}
\centering
\small
\begin{tabular}{lcccc}
\toprule
Encoder & Type & dim & $\kappa_q$ & 95\% CI \\
\midrule
CLIP-B/32 & image & 512 & 0.103 & $[.061, .140]$\\
DINOv2-B/14 & image & 768 & 0.107 & $[.072, .143]$\\
TIPSv2-B/14 & image & 768 & 0.094 & $[.058, .131]$\\
CLIP-L/14 & image & 768 & 0.189 & $[.143, .230]$\\
SigLIP-L/16 & image & 1024 & \textbf{0.247} & $[.198, .290]$\\
SigLIP-SO400M-14 & image & 1152 & 0.183 & $[.135, .226]$\\
\hline
VideoMAE-base (Kinetics) & video, 16fr & 768 & 0.101 & $[.054, .144]$\\
\bottomrule
\end{tabular}
\end{table}

Among image encoders, the scale-up from B/14 ($\sim 0.10$) to L/14
($\sim 0.20$) is the single largest gain we observed; the further
scale-up to SO400M does \emph{not} extend the trend. We attribute
the SO400M regression to val/test discrepancy under increased model
capacity on $\sim 5{,}300$ train clips: val $\kappa_q$ on SO400M is
$0.186$ versus test $\kappa_q$ of $0.183$, indicating the model is
at the edge of its supportable capacity.

\paragraph{Video-temporal pretraining does not help.} VideoMAE-base
\cite{tong2022videomae}, an obvious candidate for capturing temporal
dynamics of engagement, achieves only $\kappa_q = 0.101$ [0.054,
0.144] under the same bagged-LR + threshold recipe ---
\emph{worse than image-pre-trained CLIP-B/32 frozen features.} Two
factors plausibly explain this: (1) VideoMAE-base is pre-trained on
Kinetics action recognition, whose representational priors target
``what action is happening'' rather than ``what facial state'' --- a
distributional mismatch with classroom engagement; (2) the 16-frame
mean-pooled representation may average out the precise micro-facial
cues that engagement appears to be encoded in within image features.
The implication is sharper than ``VideoMAE-base is suboptimal'': the
\emph{structural ceiling extends across modality}. Larger video models
(VideoMAE-Large, InternVideo2) and engagement-aware video pretraining
are open directions; the current best frozen-feature path remains an
image encoder (SigLIP-L) with a careful readout.

\subsection{Reviewer-attack reruns}
\label{sec:attacks}

We rerun the experiments addressed in our prior CVPR submission's
rejection, addressing each reviewer concern.

\paragraph{Full test set (R2).} All DAiSEE results in this paper
use the full 1,784-clip subject-disjoint test split; we do not
sub-sample. (Our prior submission's $N = 300$ stratified sample
inflated $\kappa$ by roughly $+0.05$ via class-balancing artifact.)

\paragraph{GPT-4o refusals (R3).} Tab.~\ref{tab:gpt4o-refusals}
reports refusal rate alongside the metric, per Reviewer mfAP's
recommendation. P3 (chain-of-thought) on individual student faces
refuses 98.5\% of requests, making the resulting $\kappa$
uninterpretable. We drop P3 from headline tables.

\begin{table}[t]
\caption{GPT-4o on full DAiSEE test with explicit refusal rate.}
\label{tab:gpt4o-refusals}
\centering
\small
\begin{tabular}{lccc}
\toprule
Prompt & Refusals & $\kappa_q$ (kept) & $n_{\mathrm{kept}}$\\
\midrule
P1 minimal & 44 (2.5\%) & 0.021 & 1,740\\
P2 rubric & 203 (11.4\%) & 0.011 & 1,581\\
P3 CoT & 1,758 (98.5\%) & --- & 26\\
\bottomrule
\end{tabular}
\end{table}

\paragraph{Multi-frame (R1).} The single-frame vs.\ 3-frame-mean
comparison on SigLIP-L probes shows a $+0.01$ improvement
(0.213~$\rightarrow$~0.224 argmax, 0.238~$\rightarrow$~0.245 with
threshold). The improvement is real but modest. We use 3-frame mean
throughout the headline results in this paper. The reviewer
concern that the prior single-frame protocol invalidated the
temporal task is addressed; the multi-frame numbers preserve the
ceiling phenomenon.

\paragraph{Class imbalance (R4).} All probes use
\verb|class_weight='balanced'|, and the headline recipe includes
val-tuned ordinal thresholds---an interventions-style remedy that
is strictly more flexible than Menon's logit
adjustment~\cite{menon2021logit}.

\paragraph{Precision match (R6).} We rerun LLaVA-1.5-7B at full
precision via the official HuggingFace checkpoint (rather than the
4-bit Ollama quantization used in the CVPR submission). On full
DAiSEE test, $\kappa_q$ for LLaVA-1.5-7B P3 CoT shifts from $0.101$
(quantized, $N=300$) to $0.082$ (full-precision, $N=1784$). The
shift is consistent with the broader downward shift on the full
test set; LLaVA-1.5-7B does not exceed the CLIP/GPT-4o family
under either precision.

\subsection{Temporal transformer head (TEAM)}
\label{sec:team}

A natural objection to the structural-ceiling claim is that we have
not given the model proper access to temporal context. We address
this with TEAM (Temporal Engagement via Attention Modulation), a
small temporal transformer head trained end-to-end on three frames
($t = 2, 5, 8$~s) of frozen SigLIP-L features.

\textbf{Architecture.} The head projects each 1024-d frame feature
to $d_h = 256$, adds learned positional embeddings, runs a 2-layer
transformer (4 heads, FFN $4 d_h$, GELU, pre-LN, dropout 0.1), then
mean-pools over the three frames and emits three CORN ordinal
logits~\cite{shi2023corn}.

\textbf{Training.} Class-balanced weighted random sampling on
training labels; CORN binary BCE loss with per-task positive
weights derived from class frequency; AdamW
($\mathrm{lr}=10^{-3}$, weight decay $0.01$); 100-step warmup,
cosine LR decay; 25 epochs; effective batch 64. Three seeds; final
prediction is the seed-averaged sigmoid-summed expected-class
$\mathbb{E}[y]$ score, decoded via val-tuned ordinal thresholds.

\textbf{Result.} TEAM ensemble (3 seeds, 25 epochs each) on the
full DAiSEE test set:
\begin{itemize}\itemsep0pt
  \item Ensemble: $\kappa_q^{\mathrm{argmax}} = 0.084$,
   $\kappa_q^{\mathrm{thresh}} = 0.211$ [0.170, 0.258]
  \item Per-seed best (val-selected): test $\kappa_q \in \{0.216,
   0.233, 0.181\}$, mean $0.210$
  \item Strongest single epoch (oracle across all 75 epoch-seed
   pairs): test $\kappa_q = 0.254$ (seed 0, epoch 1)
\end{itemize}

\textbf{Interpretation.} The TEAM ensemble at $\kappa_q = 0.211$
sits below the frozen-feature ceiling at $\kappa_q = 0.247$ by
$-0.036$. Three observations:
\begin{itemize}\itemsep0pt
  \item \emph{Three-frame temporal modelling adds nothing over
   multi-frame averaging.} Mean-pooling SigLIP-L over $t = 2, 5, 8$
   gives $\kappa_q = 0.245$ as a linear probe. Learning a temporal
   transformer on the same three frames gives $\kappa_q = 0.211$.
   The architectural complexity does not help.
  \item \emph{Severe val/test discrepancy persists.} Val $\kappa_q$
   peaks at $0.28$ across multiple epochs while test $\kappa_q$
   oscillates between $0.15$ and $0.25$. Val-based model selection
   does not reliably pick the test peak.
  \item \emph{Individual epochs briefly exceed baseline.} Seed 0
   epoch 1 reaches test $\kappa_q = 0.254$. This is non-replicable
   across seeds (seed 42 best is $0.236$, seed 2025 best is
   $0.218$) and unselectable from val.
\end{itemize}

This is the fifth independent negative result on encoder
adaptation (after LEACE, \textsc{Dream}, VideoMAE-base, and the
SigLIP-L last-block fine-tune). The temporal-architecture lever
does not break the structural ceiling.

\subsection{End-to-end fine-tuning of a smaller backbone}
\label{sec:resnet18}

If shallow fine-tuning of a large encoder (SigLIP-L last block) fails,
a smaller backbone might be a better fit for the modest scale of
DAiSEE's training set (5,358 clips). We test this by fully fine-tuning
\textbf{ResNet18}~\cite{he2016resnet} (11.7M parameters, ImageNet
pretrained) end-to-end on DAiSEE single-frame engagement
classification. Setup: class-balanced focal loss
($\gamma = 2.0$), AdamW (encoder LR $10^{-4}$, head LR $10^{-3}$),
linear warmup + cosine decay, 15 epochs, batch 32, light augmentation
(RandomCrop, HorizontalFlip, ColorJitter), 3 seeds with ensemble
averaging.

\begin{table}[t]
\caption{ResNet18 end-to-end fine-tuning on DAiSEE. Ensemble averages
final-epoch test predictions across 3 seeds; thresholds val-tuned on
the ensemble val predictions.}
\label{tab:resnet18}
\centering
\small
\begin{tabular}{lcc}
\toprule
Configuration & test $\kappa_q$ & 95\% CI \\
\midrule
Seed 0 final epoch & 0.104 & --- \\
Seed 42 final epoch & 0.103 & --- \\
Seed 2025 final epoch & 0.137 & --- \\
Per-seed best epoch (oracle) & 0.170 & (seed 2025 ep8) \\
\textbf{3-seed ensemble + threshold} & \textbf{0.114} & $[.073, .154]$ \\
\midrule
Frozen baseline (bagged SigLIP-L + threshold) & 0.247 & --- \\
$\Delta$ (ResNet18 e2e $-$ frozen SigLIP-L) & $-0.133$ & --- \\
\bottomrule
\end{tabular}
\end{table}

ResNet18 end-to-end \emph{substantially underperforms} the frozen
SigLIP-L baseline by $-0.133$. The smaller encoder does not benefit
from the deeper adaptation. We attribute this to (a) ResNet18's
ImageNet pretraining lacks the affect-relevant inductive bias of
contrastive vision-language pretraining; (b) the modest scale of
DAiSEE's training set (5,358 clips) is insufficient to retrain a
randomly-initialized affect-discriminating representation from
scratch; and (c) the structural ceiling identified in
Sec.~\ref{sec:ceiling} appears to be a property of the engagement
\emph{signal} as captured in DAiSEE, not of any particular encoder
architecture.

\subsection{Last-block encoder fine-tuning}
\label{sec:finetune}

The strongest test of the structural-ceiling hypothesis is whether
allowing the encoder to update its weights breaks through.
We unfreeze the last transformer block + post-layernorm of
SigLIP-L/16 along with a 3-layer classification head, and train
for 6 epochs with class-balanced cross-entropy, AdamW
($\mathrm{lr}_{\mathrm{enc}}=3{\times}10^{-5}$,
$\mathrm{lr}_{\mathrm{head}}=5{\times}10^{-4}$), linear warmup,
cosine LR decay, and effective batch 64 (per-step batch 16 with
4-step gradient accumulation). Trainable parameters: 12.6M encoder
+ 1.06M head $=$ 13.7M.

\begin{table}[t]
\caption{Last-block SigLIP-L fine-tune trajectory on DAiSEE.
Val-selected best checkpoint has $\kappa_q^{\mathrm{test}}=0.110$;
oracle test peak (ep1) is $\kappa_q^{\mathrm{test}}=0.190$. Even
the oracle peak is below the frozen baseline ($\kappa_q=0.247$).}
\label{tab:finetune}
\centering
\small
\begin{tabular}{ccc}
\toprule
Epoch & val $\kappa_q$ & test $\kappa_q$ \\
\midrule
0 & 0.076 & 0.159\\
1 & 0.137 & \textbf{0.190} (oracle test peak)\\
2 & 0.177 & 0.150\\
3 & 0.177 & 0.158\\
4 & \textbf{0.207} (val peak, saved) & 0.110\\
5 & 0.190 & 0.143\\
\midrule
Frozen baseline (bagged + thresh) & --- & \textbf{0.247}\\
\bottomrule
\end{tabular}
\end{table}

Two observations:

\paragraph{Val/test divergence is severe.} By epoch~4, val
$\kappa_q = 0.207$ but test $\kappa_q = 0.110$---a gap of $+0.097$.
The model has overfit to validation-specific patterns that do not
transfer to test. Across the trajectory, the val-test correlation
is actually \emph{negative} ($\rho = -0.71$). Val-based model
selection in this regime systematically picks the worst test
epoch. The DAiSEE validation set (1{,}429 clips, 19 subjects) is
too small relative to encoder capacity to act as a reliable
selection criterion for fine-tuned models. This is itself a
methodological caution for any future DAiSEE work in this regime.

\paragraph{Even the oracle test peak does not break the ceiling.}
$\kappa_q^{\mathrm{test}} = 0.190$ at epoch~1 is the best test
$\kappa_q$ the trajectory achieves---still $-0.057$ below the
frozen bagged-LR + threshold recipe. Last-block encoder
fine-tuning does not unlock the engagement-orthogonal
representation that we hoped would expose. This is the fourth
negative data point on adapt-the-encoder paths in our study
(prior three: LEACE, \textsc{Dream}, VideoMAE-base). The
implication for the field is sharper than a single result: a
lightly fine-tuned image encoder cannot break the ceiling, a
video-temporal encoder pre-trained on Kinetics cannot break the
ceiling, and a self-supervised personalization scheme cannot
break the ceiling. The structural diagnosis stands.

\paragraph{What does break the ceiling.} The supervised SOTA on
DAiSEE (ViBED-Net, $\kappa \approx 0.55$ corresponding to 73\%
accuracy) trains a 3-D temporal CNN end-to-end on engagement
labels. End-to-end deep retraining of a video-temporal encoder on
engagement supervision \emph{does} break the ceiling. The
takeaway is that the ceiling is a property of
\emph{representation-level adaptation}, not of the task per se:
deep label-supervised training on engagement can extract a useful
representation, but shallower adaptations cannot.

\subsection{Ablations on the SOTA recipe}
\label{sec:abl-recipe}

Tab.~\ref{tab:abl-recipe} ablates each component of the recipe.

\begin{table}[t]
\caption{Ablations on the SOTA recipe (SigLIP-L). The most
impactful single intervention is val-tuned ordinal threshold
tuning. Random feature-subspace bagging (RSB) adds $\le 0.005$ over
vanilla bagging.}
\label{tab:abl-recipe}
\centering
\small
\begin{tabular}{lc}
\toprule
Recipe & $\kappa_q$ \\
\midrule
LR, argmax & 0.199\\
LR, threshold-tuned & 0.224 \;\;\;($+0.025$)\\
Bagged LR (K=20$\times$5), argmax & 0.213 \;\;($+0.014$)\\
Bagged LR (K=20$\times$5), threshold-tuned & \textbf{0.238} \;\;($+0.039$)\\
Bagged LR + class-prior calibration & 0.224\\
Bagged LR + Calibrate-Before-Use & 0.213\\
Bagged LR + RSB (frac=0.45), threshold-tuned & 0.214\\
Bagged MLP head, threshold-tuned & 0.224\\
\bottomrule
\end{tabular}
\end{table}

The single largest single-intervention is threshold tuning. Other
post-hoc methods, including those introduced for related problems
(content-free bias, class-prior adjustment), do not stack on top
of threshold tuning, because threshold tuning is itself a flexible
ordinal calibration.
